% Generated draft: requires mathematical and source review.
\documentclass[11pt]{article}
\usepackage[margin=25mm]{geometry}
\usepackage{fontspec}
\setmainfont{Latin Modern Roman}
\usepackage{amsmath,amssymb,mathtools,unicode-math}
\setmathfont{Latin Modern Math}
\usepackage{enumitem,needspace,xurl,hyperref}
\setlength{\parindent}{0pt}
\setlength{\parskip}{5pt}
\setlength{\emergencystretch}{3em}
\newcommand{\N}{\mathbb{N}}
\newcommand{\Z}{\mathbb{Z}}
\newcommand{\Q}{\mathbb{Q}}
\newcommand{\R}{\mathbb{R}}
\newcommand{\C}{\mathbb{C}}
\newcommand{\F}{\mathbb{F}}
\newcommand{\A}{\mathbb{A}}
\newcommand{\PP}{\mathbb{P}}
\newcommand{\E}{\mathbb{E}}
\newcommand{\eps}{\varepsilon}
\newcommand{\dd}{\,\mathrm{d}}
\begin{document}
\section*{UnsolvedMath ID: 20002320}
\textbf{Source problem code:} catalogue entry
\subsection{Definitions and mathematical statement}
Millar's Question 18 asks when the prime model of a decidable theory with countably many countable models must have a decidable presentation, and records that \(\omega\)-stability suffices. In the aid's setting let \(T\) be complete and decidable in a computable countable language, with \(I(T,\aleph_0)\leq\aleph_0\). A decidable presentation means that its full elementary diagram, with parameters from the presentation, is computable; this is stronger than a computable atomic diagram. A prime model embeds elementarily into every model of \(T\).

The aid reports that Millar's finite-constant-extension theorem forces some almost-prime model to have a decidable presentation. Here almost-prime means prime over a finite tuple. Order such isomorphism types by elementary embeddability, the fundamental-order convention in use. If decidability of almost-prime types is downward closed in this order, the prime model is therefore decidable. The report identifies the published positive AL-theory case through this descent property, but does not give a full AL definition or citation.

An independent sufficient condition in the aid is an effective isolation basis. A \(T\)-consistent formula \(\theta(\bar x)\) is complete if it decides every formula in those variables modulo \(T\). Suppose there is a computable enumeration of such formulas in all finite arities, with witnesses sufficiently dense that every consistent formula has an enumerated complete refinement. The aid reports that this gives a decidable prime model by an effective atomic construction. A computable list of types without effective isolation witnesses is not silently substituted for that assumption.

Consequently the reported counterexample filter is conditional: if the source hypothesis holds but the prime model has no decidable copy, then a decidable almost-prime model still exists, downward decidability descent fails, the model is outside the reported positive AL class, and no such effective principal-type/isolation basis is available. In particular the full set of complete formulas cannot be c.e. in this setting, because its enumeration together with the needed density would give the positive construction. The report's abbreviated phrase ``no c.e. set of complete formulas'' cannot literally exclude every c.e. subset, since the empty set is always c.e.; the meaningful condition concerns the full set or an adequate dense basis.

The aid also describes the hypothetical prime model as ``almost decidable,'' but does not define that technical term or state the theorem behind it. That reported qualification is retained without replacing it by a guessed degree or presentation property. The filter and sufficient conditions do not construct a counterexample or characterize every small decidable theory, so the original general question remains open in the report.
\subsection{Short English statement}
The supplied aid gives two routes to a decidable prime model: descent from a forced decidable almost-prime model, and an effective basis of isolating formulas. Any counterexample must avoid both routes and the reported AL positive class. Several abbreviated technical conditions need source clarification, so this is a filter rather than a full characterization.
\subsection{Sources}
\begin{itemize}
\item Local source record: catalogue entry, catalogue entry; canonical \texttt{aim-logic-notes.json}, \texttt{notes[95]}, Question 18 attributed to Millar. Title: ``Effective isolation and descent to the prime model.''
\item AIM workshop list \emph{Computable stability theory}, supplied reference \url{https://aimath.org/pastworkshops/computestabproblems.pdf}. The original note says \(\omega\)-stability suffices. This draft uses the supplied record without accessing a PDF.
\item Supplied attempt 1 invokes Millar's finite-constant-extension theorem, a published AL-theory result, effective isolation and an almost-decidable qualification, but supplies no full bibliographic entries or definition of the latter term.
\item Registry dated 14 August 2026: ``partially\_solved,'' ``PARTIAL-PROGRESS,'' statement status ``exact''; attempt partial, research-run status open, rigor ``proved,'' novelty confidence low. The machine-generated, non-peer-reviewed aid requires independent verification before citation.
\end{itemize}
\end{document}
